\section{Kinetic models}\label{sec:models}

\subsection{The bilinear model}

Every model in PyRATE-TA produces a concentration matrix \Cmat{} of shape
$N_\text{delays} \times N_\text{components}$. The data matrix is then bilinear,

\begin{equation}
  \Dmat \;\approx\; \Cmat\,\Smat^{\mathsf{T}},
  \label{eq:bilinear}
\end{equation}

with $\Smat$ of shape $N_\text{pixels} \times N_\text{components}$ holding the
component spectra. The kinetics enter only through \Cmat, and the spectra only
through \Smat; this separation is what makes global analysis tractable
(section~\ref{sec:varpro}).

\subsection{Model families}

\begin{table}[h]
  \centering
  \begin{tabular}{lll}
    \toprule
    \code{ModelType} & Scheme & Spectra obtained \\
    \midrule
    \code{Parallel}   & independent decays          & DAS (decay-associated) \\
    \code{Sequential} & $A \to B \to C \to \cdots$  & EAS (evolution-associated) \\
    \code{Target}     & general \Kmat, with branching & SAS (species-associated) \\
    \bottomrule
  \end{tabular}
  \caption{Model families and the meaning of the resulting spectra.}
  \label{tab:models}
\end{table}

The distinction is not cosmetic. A parallel model makes no claim about
connectivity, and its DAS are linear combinations of the true species spectra;
a sequential model assumes a single irreversible chain, and its EAS are the
spectra of the successive evolving states; only a target model with an
explicitly justified \Kmat{} yields spectra of actual species. The
\code{modelType} string handed to PyMORGAN's \code{plot\_species\_spectra}
must reflect which of the three was used, because it drives the legend wording
upstream --- mislabelling a DAS as an SAS is a scientific error, not a
formatting one.

\subsection{Parallel and sequential models}

For a parallel model with lifetimes $\tau_i$,

\begin{equation}
  C_{ni} = \exp\!\left(-t_n/\tau_i\right),
\end{equation}

and for a sequential model the same expression is replaced by the analytic
solution of the chain, which for distinct lifetimes is a sum of exponentials
with the well-known amplitude factors
$\prod_{j\neq i}\left(\tau_i^{-1}-\tau_j^{-1}\right)^{-1}$. Both are special
cases of the rate matrix below, and are implemented as such wherever that costs
nothing, so there is one propagator to test rather than three.

\subsection{Rate matrices}\label{sec:kmat}

A target model is defined by its rate matrix \Kmat, with off-diagonal element
$K_{ij}$ the rate from compartment $j$ to compartment $i$ and the diagonal the
negative sum of all rates leaving that compartment. The populations obey

\begin{equation}
  \frac{\mathrm{d}\mathbf{c}(t)}{\mathrm{d}t} = \Kmat\,\mathbf{c}(t),
  \qquad
  \mathbf{c}(0) = \mathbf{c}_0 ,
\end{equation}

so that $\mathbf{c}(t) = \exp(\Kmat t)\,\mathbf{c}_0$. \Kmat{} is built
explicitly and \Cmat{} obtained by eigendecomposition where \Kmat{} is
diagonalisable, and by a matrix-exponential propagator otherwise.

\paragraph{Exactly degenerate lifetimes.} When two rates coincide, \Kmat{} is
defective: the eigenvector matrix is singular and the eigenvalue solution does
not exist, although the kinetics are perfectly well behaved (for
$A \rightarrow B \rightarrow$ with equal rates the closed form is
$c_B = kt\,e^{-kt}$). PyRATE-TA splits the degeneracy by an antisymmetric diagonal
perturbation of relative size $10^{-6}$ and continues. This is not a cosmetic
choice: an optimiser routinely passes through equal lifetimes on its way to the
optimum, and refusing to evaluate there would abort otherwise sound fits. The
perturbation is far below any experimental precision and is logged.

\paragraph{Near-degenerate lifetimes.} When two eigenvalues approach each
other, the eigenvector matrix becomes ill-conditioned. The symptom in a fit is a
pair of amplitudes of opposite sign that grow without bound while the residual
barely improves. PyRATE-TA guards this case explicitly and warns when two fitted
lifetimes come within \code{degeneracy\_warn\_ratio} of each other
(chapter~\ref{sec:settings}); a fit that trips this warning should be re-run
with one component fewer before its lifetimes are believed.

\subsection{Predefined schemes}

Branched schemes are kept in \code{pyrate_ta.models.TARGET\_SCHEMES}, keyed by a
short identifier: equilibria (\code{A\_eq\_B} and its decaying variants),
chains with an equilibrium (\code{A\_eq\_B\_to\_C}, \code{A\_eq\_B\_eq\_C}),
branchings into a common product (\code{branch\_to\_common\_D}) and parallel
chains (\code{two\_chains\_decaying\_ends}). A scheme carries its own species
and rate counts, and these need not agree: \code{A <=> B; A ->; B ->} has two
compartments and four rate constants. A target model therefore takes its size
from the scheme, not from the requested number of components, and the number of
free lifetimes is the scheme's rate count.

\begin{lstlisting}
import pyrate_ta as pr

model = pr.make_model("Target", scheme="A_eq_B_to_C")
C = model.concentrations(t, taus=[2.0, 5.0, 20.0, 200.0], irf_fwhm=0.15)
\end{lstlisting}

\subsection{Writing a scheme}\label{sec:notation}

A target model is written as its reactions, not as a matrix:

\begin{lstlisting}
A  ->  B   : k1          # a step
B <->  C   : k2, k3      # an equilibrium (forward, backward)
C  ->      : k4          # decay to the ground state
init A = 1               # optional initial populations
\end{lstlisting}

\code{pyrate\_ta.scheme\_from\_text} turns this into a scheme, and the rate matrix
is \emph{derived}: every reaction contributes $-k$ to the reactant's diagonal
and $+k$ to the product's row, so the columns cannot fail to balance. Three
consequences are worth stating, because they are what make the notation
general:

\begin{itemize}
  \item \textbf{Species are named.} \code{S1}, \code{ICT}, \code{T1} appear on
    the diagram, in the lifetime table and in the result, so a scheme reads as
    chemistry rather than as indices.
  \item \textbf{Rates are named, so they can be shared.} The same rate name on
    two arrows is \emph{one} fitted parameter used twice --- the natural way to
    impose ``these two channels have the same rate'', which a rate matrix
    cannot express without a hand-written builder.
  \item \textbf{Any topology.} Chains, branches, equilibria, cycles,
    disconnected pools and non-decaying products are all just lines; nothing is
    special-cased, and \code{scheme\_to\_text} regenerates the notation so a
    scheme round-trips.
\end{itemize}

A malformed scheme is refused with the offending line number, never
half-built; \code{check\_scheme\_text} returns that message without raising,
which is what the editor's \emph{Update K graph} button reports.

\subsection{Scheme diagrams}

\code{pyrate_ta.plot\_scheme} draws the compartment diagram of any model, scheme
or fit result: one node per compartment, one arrow per rate, the ground state
as a separate node, and an equilibrium as two opposed arrows. Arrows are
labelled with the \emph{rate constants themselves} --- $k_1$, $k_2$, and a sum
such as $k_1+k_2$ where a compartment branches --- rather than with numbers, so
the diagram tells you which row of the rate table an arrow corresponds to when
setting values and bounds. The mapping is derived, not assumed: feeding the
scheme's $k \rightarrow \Kmat$ builder a unit vector shows exactly where each
constant appears. Pass \code{rate\_labels="value"} for the lifetime of each
step instead.

Every family is laid out \textbf{horizontally}: population flows left to right
and the ground state sits below, so a branched scheme spreads across a wide,
short figure instead of stacking a column that wastes the width and shrinks the
nodes.

The scheme is recovered from whatever the object carries. A fit records its
scheme twice --- as the notation (\code{scheme\_text}) and as a key --- and the
\emph{text} is parsed first, so a scheme written in the editor keeps its species
names and rate symbols on the diagram. This matters because a typed scheme is
stored under the key \code{"custom"}, which is in no registry: consulting the
key alone made drawing one fail outright.

\subsection{Instrument response}\label{sec:irf}

The instrument response is convolved into \Cmat, not into the data. For
exponential models this is done analytically: the convolution of a Gaussian of
width $\sigma$ centred at $t_0$ with a decaying exponential is an exponentially
modified Gaussian,

\begin{equation}
  \left(g \ast e^{-t/\tau}\right)(t)
  = \tfrac{1}{2}\exp\!\left(\frac{\sigma^{2}}{2\tau^{2}} - \frac{t-t_0}{\tau}\right)
    \operatorname{erfc}\!\left(\frac{1}{\sqrt{2}}\left(\frac{\sigma}{\tau} - \frac{t-t_0}{\sigma}\right)\right),
  \label{eq:emg}
\end{equation}

which is both faster and better conditioned than numerical convolution. The
three handling modes are \code{None} (no IRF), \code{Gaussian}
(equation~\ref{eq:emg}) and \code{Clip} (drop the delays below \code{t\_min}
instead of modelling the response).

\subsection{The coherent artefact}

The signal around $t=0$ contains contributions --- cross-phase modulation,
two-photon absorption, stimulated Raman --- that no kinetic model describes.
Two treatments are supported: clipping the early delays, or adding artefact
components (the IRF and its first derivatives) to the model. Clipping changes
which data the fit saw, so the applied \code{t\_min} is part of the result's
reproducibility payload (section~\ref{sec:repro}) and not merely a plotting
choice.

\subsection{Chirp}

Chirp is corrected \emph{upstream}, in PyMORGAN, before the data reach a fit.
PyRATE-TA does not implement a second chirp correction, and a fit of visibly
chirped data is a data-processing mistake rather than a modelling one.
